\section{Introduction}
\label{sec:intro}

Combinatorial interaction testing (CIT) is one of the most successful
systematic test-design techniques in software engineering. Its empirical
foundation is the \emph{interaction rule}: most faults are triggered by the
interaction of only a few input parameters, so a $t$-way \emph{covering array}
over the inputs---exercising every combination of values for every $t$ parameters
at least once---detects a large fraction of faults with a test suite whose size
grows only logarithmically in the number of parameters~\cite{kuhn2004software,
nie2011survey}. Mature tools such as ACTS/IPOG~\cite{lei2007ipog} and AETG~\cite{cohen1997aetg}
construct near-optimal covering arrays over inputs and are widely deployed.

\subsection{The input-side premise breaks for AI/ML and quantum systems}
CIT is fundamentally \emph{input-oriented}: it reasons about combinations of
input values and treats the system as a black box whose outputs are merely
checked by an oracle. This premise is well matched to configuration-style
software, where distinct input combinations tend to drive distinct control-flow
paths. It is far less appropriate for modern AI/ML and quantum systems:

\begin{itemize}
\item The input-to-behaviour map is \emph{opaque and stochastic}. A neural
network or a noisy quantum circuit realises a map $f: X \to \PY$ into a
distribution over structured outputs; there is no parameter-interaction
structure on the inputs that a covering array can target meaningfully.
\item Inputs \emph{collapse} onto behaviours. Enormous, high-dimensional regions
of the input space map to the same observable behaviour (the same predicted
class, the same confidence band, the same measurement-outcome regime). An input
covering array can therefore exercise many input interactions while touching
only a handful of distinct behaviours---leaving critical behaviours (e.g.\ a
fairness inconsistency, a non-monotone response, a low-fidelity quantum regime)
entirely uncovered.
\item What we want to \emph{cover} lives on the output side. The properties an
engineer cares about---calibration, robustness, fairness, monotonicity,
measurement fidelity---are all statements about the \emph{output/behaviour} of
the system, not about which inputs were combined.
\end{itemize}

\subsection{Inverting the paradigm}
We therefore invert combinatorial testing. Instead of covering input
combinations and observing whatever outputs result, \rnwise{} Output-Oriented
Testing:
\begin{enumerate}
\item partitions the structured output space $Y = Y_1 \times \cdots \times Y_q$
into semantic \emph{behaviour classes} via per-component partitioners
$\pi_j: Y_j \to W_j$, yielding an abstract output $\fhat(x) =
(\pi_1(y_1),\dots,\pi_q(y_q)) \in \Wspace$;
\item constructs a $t$-way \emph{Output Covering Array} over $\Wspace$, so that
every feasible $s$-way combination of output behaviours is a coverage target;
\item recovers, for each target, a concrete input by \emph{inverse mapping}---a
gradient-free search that minimises the mismatch between the realised abstract
output and the target;
\item prioritises the resulting tests by incremental output coverage.
\end{enumerate}
Coverage is measured on the output side (the fraction of feasible $s$-way output
interactions exercised), which directly reflects behavioural adequacy and is
model-agnostic: the identical procedure applies to a gradient-boosted classifier,
a neural network, a text classifier, and a quantum circuit.

\subsection{Contributions}
\begin{enumerate}
\item \textbf{A new adequacy criterion and formalism.} We define the abstract
output map $\fhat$, the Output Covering Array $\OCA(M;s,q,w)$, and output
coverage $\OCov_s$, and we establish coverage bounds---including the lower bound
$M \ge v^{s}$ (Corollary~\ref{cor:lb})---for output-side combinatorial coverage
(Section~\ref{sec:framework}).
\item \textbf{A constructive algorithm.} A four-stage procedure
(feasibility, covering-array construction, gradient-free inverse mapping,
prioritisation) that produces an executable, prioritised test suite realising the
target output interactions (Algorithm~\ref{alg:rnwise}).
\item \textbf{A domain-agnostic instantiation.} We instantiate the method with
five inverse-mapping optimisers---Jaya, Whale, Firefly, Bayesian optimisation,
and a variational (VQE-style) optimiser for quantum systems---and four benchmark
domains spanning tabular ML, vision, NLP, and quantum computing.
\item \textbf{A rigorous empirical evaluation.} Over 30 repeated runs with
non-parametric significance testing (Mann--Whitney $U$, bootstrap confidence
intervals, Cliff's $\delta$), a strength study ($s=2,3,4$), ablations, and a cost
profile, we show \rnwise{} uniquely attains perfect output coverage, matches the
strongest behavioural baseline on fault detection, and generalises across all
four domains (Section~\ref{sec:eval}).
\end{enumerate}
